
\documentclass{article}
\usepackage{colortbl}
\usepackage{makecell}
\usepackage{multirow}
\usepackage{supertabular}

\begin{document}

\newcounter{utterance}

\centering \large Interaction Transcript for game `referencegame', experiment `line\_grids\_columns', episode 16 with reflexion{-}Qwen3.5{-}4B.
\vspace{24pt}

{ \footnotesize  \setcounter{utterance}{1}
\setlength{\tabcolsep}{0pt}
\begin{supertabular}{c@{$\;$}|p{.15\linewidth}@{}p{.15\linewidth}p{.15\linewidth}p{.15\linewidth}p{.15\linewidth}p{.15\linewidth}}
    \# & $\;$A & \multicolumn{4}{c}{Game Master} & $\;\:$B\\
    \hline

    \theutterance \stepcounter{utterance}  
    & & \multicolumn{4}{p{0.6\linewidth}}{
        \cellcolor[rgb]{0.9,0.9,0.9}{
            \makecell[{{p{\linewidth}}}]{
                \texttt{\tiny{[P1$\langle$GM]}}
                \texttt{You have attempted this same task before and did not succeed. These are the plans you wrote after those attempts.} \\
\texttt{Trial \#1: New plan: My previous description was too vague because it listed the empty columns without explicitly defining the filled column's position relative to the grid's structure, which likely failed to uniquely distinguish the target from the distractors. The target grid has a solid vertical bar of 'X's in the 3rd column, while the distractors have blocks of 'X's in the 3rd/4th/5th columns (top{-}right) or 4th/5th columns (top{-}right). I will revise the expression to specifically state "a vertical column of filled cells in the third column" and emphasize that the first two columns are entirely empty, ensuring the description matches the unique pattern of the target grid (3rd column full, others empty) rather than just listing column contents.} \\
\\ 
\texttt{Trial \#2: New plan: My previous description failed because it did not explicitly state the row configuration, which is the primary difference between the target grid and the distractors. The target grid has a solid horizontal row of filled cells in the third row, whereas the distractors have filled cells only in the top rows. I will revise the expression to specifically state "a solid horizontal row of filled cells in the third row" and clarify that all other cells are empty, ensuring the description uniquely identifies the target grid based on its row structure rather than just column properties.} \\
\\ 
\texttt{Trial \#3: New plan: My previous attempts failed because I misidentified the primary structural feature of the target grid. I incorrectly assumed the target had a "solid vertical column" (Trial \#1) or a "solid horizontal row" (Trial \#2), but the actual target grid consists of a **solid vertical column of filled cells in the third column** where the first two columns are entirely empty. The distractors differ by having filled cells in the fourth and fifth columns (Distractor 1) or only in the fourth and fifth columns (Distractor 2). To succeed, I must generate an expression that explicitly states the target grid has a **full third column** while the **first two columns are completely empty**, thereby distinguishing it from the distractors which lack filled cells in the first two columns and have filled cells in the rightmost columns. The expression must be: "Expression: The grid has a solid vertical column of filled cells in the third column, with the first two columns entirely empty."} \\
\\ 
\\ 
\texttt{You are given three grids, where each of them is 5 by 5 in size.} \\
\texttt{Grids have empty cells marked with "▢" and filled cells marked with "X".} \\
\texttt{Your task is to generate a referring expression that best describes the target grid while distinguishing it from the two other distractor grids.} \\
\texttt{The first grid is the target grid, and the following two grids are the distractors.} \\
\\ 
\texttt{Target grid:} \\
\\ 
\texttt{X □ X X □} \\
\texttt{X □ X X □} \\
\texttt{X □ X X □} \\
\texttt{X □ X X □} \\
\texttt{X □ X X □} \\
\\ 
\texttt{Distractor grid 1:} \\
\\ 
\texttt{□ □ X X X} \\
\texttt{□ □ X X X} \\
\texttt{□ □ X X X} \\
\texttt{□ □ X X X} \\
\texttt{□ □ X X X} \\
\\ 
\texttt{Distractor grid 2:} \\
\\ 
\texttt{X □ □ X X} \\
\texttt{X □ □ X X} \\
\texttt{X □ □ X X} \\
\texttt{X □ □ X X} \\
\texttt{X □ □ X X} \\
\\ 
\texttt{Instruction: Describe the target grid.} \\
\texttt{Generate the referring expression starting with the tag "Expression: " for the given target grid. Omit any other text.} \\
            }
        }
    }
    & \\ \\

    \theutterance \stepcounter{utterance}  
    & \multicolumn{4}{p{0.6\linewidth}}{
        \cellcolor[rgb]{0.9,0.9,0.9}{
            \makecell[{{p{\linewidth}}}]{
                \texttt{\tiny{[P1$\rangle$GM]}}
                \texttt{Expression: The grid has a solid vertical column of filled cells in the third column, with the first two columns entirely empty.} \\
            }
        }
    }
    & & \\ \\

    \theutterance \stepcounter{utterance}  
    & & & \multicolumn{2}{p{0.3\linewidth}}{
        \cellcolor[rgb]{0.9,0.9,0.9}{
            \makecell[{{p{\linewidth}}}]{
                \texttt{\tiny{[GM$|$GM]}}
                \texttt{Expression: The grid has a solid vertical column of filled cells in the third column, with the first two columns entirely empty.} \\
            }
        }
    }
    & & \\ \\

    \theutterance \stepcounter{utterance}  
    & & \multicolumn{4}{p{0.6\linewidth}}{
        \cellcolor[rgb]{0.9,0.9,0.9}{
            \makecell[{{p{\linewidth}}}]{
                \texttt{\tiny{[GM$\rangle$P2]}}
                \texttt{You have attempted this same task before and did not succeed. These are the plans you wrote after those attempts.} \\
\texttt{Trial \#1: New plan: I incorrectly identified the target grid by focusing on the filled cells in the third column but failing to verify the status of the other columns. The expression requires the first, second, fourth, and fifth columns to be entirely empty. In the "First" grid, the first two columns are indeed empty, but I did not explicitly confirm that the fourth and fifth columns are also empty before selecting. Next time, I will systematically check every column of the candidate grid against the expression's requirements: confirm the specified column is filled and all other columns are completely empty before answering.} \\
\\ 
\texttt{Trial \#2: New plan: I incorrectly identified the target grid by focusing on the filled cells in the third column but failing to verify the status of the other columns. The expression requires the first, second, fourth, and fifth columns to be entirely empty. In the "First" grid, the first two columns are indeed empty, but I did not explicitly confirm that the fourth and fifth columns are also empty before selecting. Next time, I will systematically check every column of the candidate grid against the expression's requirements: confirm the specified column is filled and all other columns are completely empty before answering.} \\
\\ 
\texttt{Trial \#3: New plan: I incorrectly identified the target grid by assuming the "First" grid matched the description because its third column was filled, but I failed to verify that the fourth and fifth columns were actually empty. In the "First" grid, the fourth and fifth columns contain filled cells ("X"), which violates the requirement that "all other cells are empty." Next time, I will strictly verify that the specified column is fully filled AND that every single cell in all other columns is empty before selecting a grid. I will compare the "Second" and "Third" grids against this full criteria to ensure no partial matches are chosen.} \\
\\ 
\\ 
\texttt{You are given three grids, where each of them is 5 by 5 in size.} \\
\texttt{Grids have empty cells marked with "▢" and filled cells marked with "X".} \\
\texttt{You are also given a referring expression that describes one of the given grids.} \\
\texttt{Your task is to select the grid that matches the given referring expression.} \\
\texttt{Generate only the number (in text) of the grid that matches the given expression by selecting first, second, or third.} \\
\\ 
\texttt{First:} \\
\\ 
\texttt{□ □ X X X} \\
\texttt{□ □ X X X} \\
\texttt{□ □ X X X} \\
\texttt{□ □ X X X} \\
\texttt{□ □ X X X} \\
\\ 
\texttt{Second:} \\
\\ 
\texttt{X □ X X □} \\
\texttt{X □ X X □} \\
\texttt{X □ X X □} \\
\texttt{X □ X X □} \\
\texttt{X □ X X □} \\
\\ 
\texttt{Third:} \\
\\ 
\texttt{X □ □ X X} \\
\texttt{X □ □ X X} \\
\texttt{X □ □ X X} \\
\texttt{X □ □ X X} \\
\texttt{X □ □ X X} \\
\\ 
\texttt{Expression: The grid has a solid vertical column of filled cells in the third column, with the first two columns entirely empty.} \\
\texttt{Question: Which grid does the expression refer to?} \\
\texttt{Start with the tag "Answer: ", followed by your selection. Omit any other text.} \\
            }
        }
    }
    & \\ \\

    \theutterance \stepcounter{utterance}  
    & & & \multicolumn{4}{p{0.6\linewidth}}{
        \cellcolor[rgb]{0.9,0.9,0.9}{
            \makecell[{{p{\linewidth}}}]{
                \texttt{\tiny{[GM$\langle$P2]}}
                \texttt{Answer: 1} \\
            }
        }
    }
     \\ \\

    \theutterance \stepcounter{utterance}  
    & & & \multicolumn{2}{p{0.3\linewidth}}{
        \cellcolor[rgb]{0.9,0.9,0.9}{
            \makecell[{{p{\linewidth}}}]{
                \texttt{\tiny{[GM$|$GM]}}
                \texttt{1} \\
            }
        }
    }
    & & \\ \\

    \theutterance \stepcounter{utterance}  
    & \multicolumn{4}{p{0.6\linewidth}}{
        \cellcolor[rgb]{0.95,0.95,0.95}{
            \makecell[{{p{\linewidth}}}]{
                \texttt{\tiny{[P2$\rightarrow$P1]}}
                \texttt{Post{-}episode Reflexion (not sent to GM):} \\
\texttt{New plan: The previous attempt failed because the description was factually correct regarding the target grid but insufficient to distinguish it from the distractors. The target grid has filled cells only in the third column, whereas Distractor 1 has filled cells in the third, fourth, and fifth columns, and Distractor 2 has filled cells only in the fourth and fifth columns. The key distinguishing feature is that the target grid has **no filled cells in the fourth or fifth columns**, while the distractors do. My new plan is to explicitly state that the grid has a solid vertical column of filled cells in the third column, and crucially, that the **fourth and fifth columns are entirely empty**, which uniquely separates it from the distractors that contain filled cells in those rightmost columns.} \\
\\ 
\texttt{Expression: The grid has a solid vertical column of filled cells in the third column, with the fourth and fifth columns entirely empty.} \\
            }
        }
    }
    & & \\ \\

    \theutterance \stepcounter{utterance}  
    & & & \multicolumn{4}{p{0.6\linewidth}}{
        \cellcolor[rgb]{0.95,0.95,0.95}{
            \makecell[{{p{\linewidth}}}]{
                \texttt{\tiny{[P1$\rightarrow$P2]}}
                \texttt{Post{-}episode Reflexion (not sent to GM):} \\
\texttt{New plan: I incorrectly selected the "First" grid because I only verified that the third column was filled and the first two columns were empty, but I failed to check the fourth and fifth columns. The expression requires the first two columns to be empty AND the third to be filled, implying the others are empty, but the "First" grid has filled cells in the fourth and fifth columns. Next time, I will strictly verify that the third column is entirely filled with "X" AND that the first, second, fourth, and fifth columns are entirely empty (containing only "▢") before selecting. I will compare the "Second" grid (which has filled cells in the first column) and the "Third" grid (which has filled cells in the first column) against this full criteria to confirm the "First" grid is the only one that fails the empty column check, ensuring I select the correct grid based on all constraints.} \\
            }
        }
    }
     \\ \\

\end{supertabular}
}

\end{document}
